\subsection{Resolutions and Grothendieck groups}

If \(\mathcal{C}\) is a set of classes of A-modules, we shall say that a left resolution \((\mathbf{P}, p)\) is bounded of type \(\mathcal{C}\) if the complex \(\mathbf{P}\) is bounded of type \(\mathcal{C}\) (X, p.~41).

THEOREM 1. --- Let \(\mathcal{C}_{0}\) and \(\mathcal{C}\) be two additive and left exact classes of A-modules such that \(\mathcal{C}_{0} \subset \mathcal{C}\) and that every A-module of type \(\mathcal{C}\) has a bounded left resolution of type \(\mathcal{C}_{0}\) . Then the homomorphism \(\alpha : \mathrm{K}(\mathcal{C}_{0}) \to \mathrm{K}(\mathcal{C})\) deduced from the inclusion of \(\mathcal{C}_{0}\) into \(\mathcal{C}\) is bijective; if M is an A-module of type \(\mathcal{C}\) and P a bounded left resolution of M of type \(\mathcal{C}_{0}\) , we have \(\alpha^{-1}([M]_{\mathcal{C}}) = \chi_{\mathcal{C}_{0}}(P)\) (X, p.~41, example 6).

Lemma 4. --- Let \(f: \mathbf{M}' \to \mathbf{M}\) be a homomorphism of A-modules of type \(\mathcal{C}\) , and \(p: \mathbf{P} \to \mathbf{M}\) a left resolution of \(\mathbf{P}\) bounded of type \(\mathcal{C}_0\) There exists a left resolution \(p': \mathbf{P}' \to \mathbf{M}'\) bounded of type \(\mathcal{C}_0\) and a morphism of complexes \(u: \mathbf{P}' \to \mathbf{P}\) such that \(p \circ u = f \circ p'\) .

Let us reason by induction on the length n of P, the assertion being trivial when this length is \textless{} 0. Consider the mapping \(g: \mathbf{M}' \times \mathbf{P}_0 \to \mathbf{M}\) such that

\[
g (x, y) = f (x) - p _ {0} (y) \quad \text { for } \quad x \in \mathbf {M} ^ {\prime}, y \in \mathbf {P} _ {0},
\]

and its kernel K; the A-module K is of type C since g is surjective and \(M' \times P_{0}\) and M are of type C. Let \(h : P_{0}' \to K\) be a surjective homomorphism where \(P_{0}'\) is of type \(C_{0}\) ; denote by \(p_{0}' : P_{0}' \to M'\) (resp. \(u_{0} : P_{0}' \to P_{0}\) ) the composite homomorphism of h and the projection \(K \to M\) (resp. \(K \to P_{0}\) ); the homomorphism \(p_{0}'\) is surjective and we have a commutative diagram

\[
\begin{array}{c} \mathbf {P} _ {0} ^ {\prime} \xrightarrow {u _ {0}} \mathbf {P} _ {0} \\ p _ {0} ^ {\prime} \Big \downarrow \qquad \qquad \qquad \qquad \qquad \qquad \qquad \qquad \qquad \Big \downarrow p _ {0} \\ \mathbf {M} ^ {\prime} \xrightarrow [ f ]{} \mathbf {M}. \end{array}
\]

It then suffices to apply the induction hypothesis to the homomorphism

\[
\operatorname{Ker} p _ {0} ^ {\prime} \rightarrow \operatorname{Ker} p _ {0}
\]

deduced from \(u_{0}\) .

Lemma 5. --- Consider a commutative diagram

\[
\begin{array}{c} \mathbf {P} ^ {\prime} \xrightarrow {u} \mathbf {P} \\ p ^ {\prime} \Big \downarrow \quad \Big \downarrow p \\ 0 \longrightarrow \mathbf {M} ^ {\prime} \xrightarrow {f} \mathbf {M} \longrightarrow \mathbf {M} ^ {\prime \prime} \longrightarrow 0 \end{array}
\]

where (P, p) (resp. (P′, p′)) is a left resolution of M (resp. M′), and where the bottom horizontal row is an exact sequence. There exists a homologism \(p''\) : Con \((u) \to \mathbf{M}''\) .

Indeed, the exact sequence (X, p.~37, prop. 7)

\[
0 \to \mathrm{P} \xrightarrow {\pi} \operatorname{Con} (u) \xrightarrow {\delta} \mathrm{P} ^ {\prime} (- 1) \to 0
\]

gives an exact homology sequence

\[
\begin{array}{c}\rightarrow \mathrm{H} _ {n} (\mathrm{P}) \rightarrow \mathrm{H} _ {n} (\mathrm{Con} (u)) \rightarrow \mathrm{H} _ {n - 1} (\mathrm{P} ^ {\prime}) \rightarrow \dots\\\dots \rightarrow \mathrm{H} _ {1} (\mathrm{Con} (u)) \rightarrow \mathrm{H} _ {0} (\mathrm{P} ^ {\prime}) \stackrel {{\partial}} {{\rightarrow}} \mathrm{H} _ {0} (\mathrm{P}) \rightarrow \mathrm{H} _ {0} (\mathrm{Con} (u)) \rightarrow 0.\end{array}
\]

According to X, p.~38, lemma 3 a), we have \(\partial = -\mathrm{H}_{0}(u)\) . Since \(\mathrm{H}_{n}(\mathbf{P}) = 0 = \mathrm{H}_{n}(\mathbf{P}')\) for \(n > 0\) and \(\mathrm{H}_{0}(u): \mathrm{H}_{0}(\mathrm{P}') \to \mathrm{H}_{0}(\mathbf{P})\) is identified with \(f: \mathbf{M}' \to \mathbf{M}\) , we conclude that \(\mathrm{H}_{n}(\mathrm{Con}(u)) = 0\) for \(n > 0\) and that \(\mathrm{H}_{0}(\mathrm{Con}(u))\) is isomorphic to \(\mathbf{M}''\) , whence the lemma.

Now let us prove the theorem.

\begin{enumerate}
\def\labelenumi{\alph{enumi})}
\tightlist
\item
  Let M be an A-module of type C. For every left resolution \((\mathbf{P}, p)\) of M, bounded of type \(C_{0}\) , the element \(\chi_{\mathcal{C}_{0}}(\mathbf{P})\) of \(\mathbf{K}(\mathcal{C}_{0})\) depends only on M. Indeed, let \((\mathbf{P}_{1}, p_{1})\) and \((\mathbf{P}_{2}, p_{2})\) be two resolutions of this type. Consider the resolution
\end{enumerate}

\[
(\mathbf {P} _ {1} \times \mathbf {P} _ {2}, p _ {1} \times p _ {2})
\]

of the A-module \(\mathbf{M} \times \mathbf{M}\) and the homomorphism \(\Delta : x \mapsto (x, x)\) of \(\mathbf{M}\) to \(\mathbf{M} \times \mathbf{M}\) . By Lemma 4, there exists a resolution \((\mathbf{Q}, q)\) of \(\mathbf{M}\) bounded of type \(\mathcal{C}_0\) and a commutative diagram

\[
\begin{array}{c c c} \mathbf {Q} & \stackrel {{u}} {{\longrightarrow}} \mathbf {P} _ {1} & \times \mathbf {P} _ {2} \\ q \Big \downarrow & & \Big \downarrow p _ {1} \times p _ {2} \\ \mathbf {M} & \stackrel {{\Delta .}} {{\longrightarrow}} \mathbf {M} & \times \mathbf {M}; \end{array}
\]

we deduce a commutative diagram

\[
\begin{array}{c c c} \mathrm{Q} & \xrightarrow {u \circ p r _ {i}} & \mathrm{P} _ {i} \\ q \Big \downarrow & & \Big \downarrow p _ {i} \\ \mathbf {M} & \xrightarrow {\mathrm{l} _ {\mathrm{M}}} & \mathbf {M}, \end{array} \qquad i = 1, 2.
\]

By lemma 5, Con \((u \circ pr_i)\) has zero homology, hence \(u \circ pr_i\) is a homologism and \(\chi_{\mathcal{C}_0}(\mathbf{Q}) = \chi_{\mathcal{C}_0}(\mathbf{P}_i)\) (X, p.~41, prop. 10); it follows that \(\chi_{\mathcal{C}_0}(\mathbf{P}_1) = \chi_{\mathcal{C}_0}(\mathbf{P}_2)\) as announced.

\begin{enumerate}
\def\labelenumi{\alph{enumi})}
\setcounter{enumi}{1}
\tightlist
\item
  For every A-module M of type C, let \(\varphi(M) \in K(\mathcal{C}_{0})\) be the common value of \(\chi_{\mathcal{C}_{0}}(P)\) for all bounded left resolutions P of M of type \(C_{0}\) . Let us show that the function \(\varphi : \mathcal{C} \to K(\mathcal{C}_{0})\) is additive. Let therefore
\end{enumerate}

\[
0 \to \mathbf {M} ^ {\prime} \xrightarrow {f} \mathbf {M} \to \mathbf {M} ^ {\prime \prime} \to 0
\]

be an exact sequence of A-modules of type C. By Lemma 4, there exists a commutative diagram

\[
\begin{array}{c} \mathrm{P} ^ {\prime} \xrightarrow {u} \mathrm{P} \\ p ^ {\prime} \Big \downarrow \qquad \qquad \qquad \Big \downarrow p \\ 0 \longrightarrow \mathrm{M} ^ {\prime} \xrightarrow {f} \mathrm{M} \longrightarrow \mathrm{M} ^ {\prime \prime} \longrightarrow 0 \end{array}
\]

where (P, p) and (P', p') are bounded left resolutions of type \(\mathcal{C}_0\) . Then we have

\[
\varphi (\mathbf {M}) = \chi_ {\mathcal {C} _ {0}} (\mathbf {P}), \quad \varphi (\mathbf {M} ^ {\prime}) = \chi_ {\mathcal {C} _ {0}} (\mathbf {P} ^ {\prime})
\]

and by Lemma 5

\[
\varphi (\mathbf {M} ^ {\prime \prime}) = \chi_ {\mathcal {C} _ {0}} (\operatorname{Con} (u)) = \chi_ {\mathcal {C} _ {0}} (\mathrm{P}) - \chi_ {\mathcal {C} _ {0}} (\mathrm{P} ^ {\prime}) = \varphi (\mathrm{M}) - \varphi (\mathrm{M} ^ {\prime});
\]

which is what we wanted to prove.

\begin{enumerate}
\def\labelenumi{\alph{enumi})}
\setcounter{enumi}{2}
\tightlist
\item
  Let then \(\beta: \mathrm{K}(\mathcal{C}) \to \mathrm{K}(\mathcal{C}_0)\) be the homomorphism such that, with the preceding notation, we have \(\beta([\mathrm{M}]_{\mathcal{C}}) = \chi_{\mathcal{C}_0}(\mathrm{P})\) . Since \(p\) is a homologism, we have \(\chi_{\mathcal{C}}(\mathrm{P}) = [\mathrm{M}]_{\mathcal{C}}\) , hence \(\alpha \circ \beta([\mathrm{M}]_{\mathcal{C}}) = \alpha(\chi_{\mathcal{C}_0}(\mathrm{P})) = \chi_{\mathcal{C}}(\mathrm{P}) = [\mathrm{M}]_{\mathcal{C}}\) and \(\alpha \circ \beta = 1_{\mathrm{K}(\mathcal{C})}\) . If M is of type \(\mathcal{C}_0\) , then (M, \(1_{\mathrm{M}}\) ) is a resolution of M, hence \(\varphi(\mathrm{M}) = [\mathrm{M}]_{\mathcal{C}_0}\) and \(\beta \circ \alpha = 1_{\mathrm{K}(\mathcal{C}_0)}\) , which completes the proof.
\end{enumerate}

We shall apply this theorem to modules of "finite projective dimension" in § 8 (X, p.~137).

